\worldchapter{Die Kunst der Synthese und des Brauens}{brewing}
\index{Synthese}\index{Brauen}
\label{ch:brewing}

\begin{multicols}{2}

Das Mischen von Chemikalien, das Sieden von Tinkturen und das Herstellen experimenteller Seren ist ein exaktes und oft gefährliches Handwerk. Um diesen Vorgang am Spieltisch dynamisch abzubilden, wird ein strukturierter Ablauf genutzt, bei dem Umstände, Ausrüstung und fehlende Komponenten direkten Einfluss auf den Herstellungswurf haben.

Das Vorgehen wird in folgenden Schritten ermittelt:
\begin{itemize}
    \item \textbf{Boni und Mali ermitteln}: Fertigkeit, Schwierigkeitsstufe der Formel und gegebenenfalls Modifikatoren durch Umgebung, Ausrüstung und ersetzte Komponenten bestimmen.
    \item \textbf{Herstellungswurf}: Mit dem errechneten Würfelpool gegen den ermittelten Mindestwurf würfeln.
    \item \textbf{Erfolge auswerten}: Prüfen, ob die benötigte Anzahl an Erfolgen erreicht wurde.
    \item \textbf{Qualität oder Patzer}: Optional hohe Qualität, Fehlschläge und gefährliche Nebenwirkungen auswerten.
\end{itemize}

\section{Die Fertigkeit}
\index{Fertigkeit!Brauen}

Das Herstellen von Substanzen basiert auf der Fertigkeit, die zum Setting und zur konkreten Substanz passt. In der Regel wird auf \textit{Chemie} oder \textit{Alchemie} gewürfelt.

In Fantasy-Settings kann eine Formel häufig mit \textit{Alchemie} oder \textit{Naturkunde} hergestellt werden. In modernen oder futuristischen Settings kommen meistens \textit{Chemie}, \textit{Naturkunde} oder \textit{Medizin} in Frage.

Wer über die primär geforderte Spezialkenntnis nicht verfügt, kann auf eine verwandte Fertigkeit ausweichen, sofern die Spielleitung diese zulässt. Der Charakter erhält dann einen Malus von \textbf{-2 Würfeln}. Die Spielleitung entscheidet, ob eine alternative Fertigkeit logisch anwendbar ist.

\begin{pexample}[Beispiel]
Der Straßendoc Patch besitzt keine Fertigkeit in \textit{Chemie} oder \textit{Naturkunde}, hat aber einen Wert von 5 in \textit{Medizin}. Er versucht, aus Haushaltsmitteln eine provisorische Rauchbombe herzustellen, um vor einem Konzerntrupp zu fliehen.

Die Spielleitung lässt \textit{Medizin} zu, da Patch rudimentäre Kenntnisse chemischer Reaktionen besitzt. Da er fachfremd improvisiert, erhält er einen Malus von \textbf{-2 Würfeln}. Statt mit 5 Würfeln wirft Patch für diesen Herstellungswurf also nur mit 3 Würfeln.
\end{pexample}

\section{Formeln und Rezepte}
\index{Formeln}\index{Rezepte}

Jedes Vorhaben setzt eine Anleitung voraus: ein Rezept, eine chemische Formel, einen Bauplan oder vergleichbare Notizen. Werden die Vorgaben nicht eingehalten, sind Misserfolge möglich. Alle Formeln sind in Schwierigkeitsgrade eingeteilt, die den Mindestwurf und die benötigten Erfolge bestimmen.

Da der übliche Mindestwurf in \trans{book_title} bei 5+ liegt, erfordern höhere Stufen oft explodierende Würfel und kritische Erfolge.

\begin{center}
\begin{tabular}{@{}lll@{}}
\toprule
\textbf{Stufe} & \textbf{MW} & \textbf{Erfolge}\\
\midrule
Einfach & 5+ & 1\\
Mittel & 8+ & 2\\
Schwer & 11+ & 4\\
\bottomrule
\end{tabular}
\end{center}

\begin{pexample}[Beispiel]
Die Agentin Bettina versucht sich an der Synthese eines schweren Wahrheitsserums. Die Tabelle gibt dafür einen Mindestwurf von 11+ vor und verlangt 4 Erfolge.

Bettina würfelt mit 5 Würfeln. Mit normalen Würfen kann sie maximal eine 6 erzielen. Sie ist also darauf angewiesen, dass mehrere ihrer Würfel explodieren, um Ergebnisse von 11 oder höher zu erreichen und die benötigten Erfolge zu sammeln.
\end{pexample}

\section{Ausrüstung und Umgebung}
\index{Ausrüstung und Umgebung}

Ein Experte wird auch durch seine Umgebung und sein Werkzeug definiert. Ausrüstung verändert den Mindestwurf und teilweise den Würfelpool.

\begin{itemize}
    \item \textbf{High-End Labor (-3 MW)}: Modernste Analytik, perfekte Werkzeuge oder ein magischer Nexus. Hier gelingt das Unmögliche.
    \item \textbf{Standard-Labor (-2 MW)}: Solide, stationäre Grundausstattung, etwa Tiegel, Zentrifugen oder Abzüge.
    \item \textbf{Mobiles Feld-Kit (-1 MW)}: Kompaktkoffer für Agenten, Feldforscher oder Abenteurer.
    \item \textbf{Basisausrüstung (+/-0 MW)}: Mörser, kleiner Kocher oder einfache Reagenzgläser. Das absolute Minimum.
    \item \textbf{Improvisiert (+1 MW, -2 Würfel)}: Behelfsgläser, hohle Rohre, bloße Hände. Ein sicheres Rezept für ein Desaster.
\end{itemize}

Die Umgebung verändert den Würfelpool:
\begin{itemize}
    \item \textbf{Labor-Idylle (+1 Würfel)}: Absolute Ruhe, kein Zeitdruck.
    \item \textbf{Normaler Alltag (+/-0 Würfel)}: Sauberer Arbeitsplatz, stabiler Tisch.
    \item \textbf{Widrige Umstände (-1 Würfel)}: Knietief im Schlamm, im fahrenden Fluchtauto oder in Eile.
    \item \textbf{Chaos (-2 Würfel)}: Aktiver Schusswechsel, einstürzende Ruine oder extremer Zeitdruck.
\end{itemize}

\begin{pexample}[Beispiel]
Toria versucht hastig, ein experimentelles Gegengift zu mischen, während Feinde bereits die verbarrikadierte Tür einschlagen. Sie nutzt ein mobiles Feld-Kit und senkt den Mindestwurf um 1.

Die Spielleitung wertet die Situation jedoch als \textit{Chaos} und zieht 2 Würfel von Torias Würfelpool ab. Ein eigentlich mittelschweres Rezept (MW 8+, 2 Erfolge) wird so mit einem MW von 7+ und einem um 2 Würfel reduzierten Pool gewürfelt.
\end{pexample}

\section{Substitution}
\index{Substitution}

In der Wildnis oder auf der Flucht ist man selten mit perfekt sortiertem Equipment gesegnet. Wenn eine benötigte Komponente fehlt, kann der Charakter nach einer Alternative suchen.

Jede Komponente in einer Formel erfüllt eine bestimmte Funktion und besitzt einen oder mehrere inhärente Effekte. Typische Effekte sind zum Beispiel Reinigung, Bindung, Erhitzung oder Stabilisierung. Anhand dieser Effekte entscheiden Spielleitung und Spieler, wie gut sich eine alternative Zutat in den Prozess einfügt.

\begin{itemize}
    \item \textbf{Passender Ersatz (+1 MW)}: Die Effekte und Eigenschaften der Stoffe decken sich gut. Der Prozess bleibt stabil, der Mindestwurf erhöht sich lediglich um einen Punkt.
    \item \textbf{Riskanter Ersatz (+2 MW)}: Die Alternative erfüllt grob denselben Zweck, ist aber flüchtig, unberechenbar oder besitzt störende Zusatzeffekte. Das sensible Gleichgewicht gerät in Gefahr.
\end{itemize}

\subsection{Das Gesetz der Instabilität}
\index{Gesetz der Instabilität}

Sobald in einem einzigen Rezept drei oder mehr Zutaten substituiert werden, wird die Substanz unberechenbar. Bei einem Misserfolg unter dieser Bedingung gilt die Probe nicht mehr als einfaches Scheitern. Die Spielleitung würfelt automatisch auf der Tabelle für Fehlschläge und Katastrophen.

\begin{pexample}[Beispiel]
Der Alchemist Kael benötigt für einen Heiltrank eigentlich süße Mondbeeren, da diese den Braueffekt \textit{Bindung} besitzen und toxische Beikräuter neutralisieren. Da er tief in den Katakomben feststeckt, greift er auf säuerliche Höhlenpilze zurück.

Die Pilze besitzen ebenfalls den Effekt \textit{Bindung}, bringen aber eine unberechenbare Eigenschaft mit sich. Die Spielleitung wertet dies als riskanten Ersatz. Kaels Mindestwurf steigt um 2 Punkte. Hätte er insgesamt drei Zutaten ersetzt, hätte das Gesetz der Instabilität gegriffen.
\end{pexample}

\section{Qualität und Patzer}
\index{Qualität}\index{Patzer}

Das Herstellen von Substanzen ist kein steriles Handwerk. Nuancen entscheiden über Meisterwerke, nutzlose Schlämme oder tödliche Gifte.

\subsection{Hohe Qualität}
\index{Hohe Qualität}

Wurden mindestens doppelt so viele Erfolge erzielt, wie die Schwierigkeitsstufe verlangt, bestimmt ein W6-Wurf die besondere Eigenschaft des Meisterstücks:

\begin{itemize}
    \item \textbf{1-2, Verstärkung}: Die Wirkung hält doppelt so lange an oder ist um 50\% effektiver.
    \item \textbf{3-4, Reinheit}: Die Substanz hat keine Nebenwirkungen und belastet den Organismus nicht.
    \item \textbf{5, Effizienz}: Die Wirkung tritt sofort und ohne Verzögerung ein.
    \item \textbf{6, Meisterstück}: Maximale Wirkung. Der Konsument erhält zusätzlich einen dauerhaften Bonus von \textbf{+1 Würfel} auf seinen nächsten relevanten Wurf.
\end{itemize}

\subsection{Fehlschlag und Katastrophe}
\index{Fehlschlag und Katastrophe}

Zeigt der Wurf mehr Einsen als tatsächliche Erfolge, gilt die Probe als gepatzt, selbst wenn die benötigten Erfolge rechnerisch erreicht wurden. Die Spielleitung würfelt 1W6 für das Ausmaß des Fehlschlags:

\begin{itemize}
    \item \textbf{1-2, Instabilität}: Die chemische oder mystische Bindung löst sich. Die Substanz verliert spätestens nach 24 Stunden ihre Wirkung und wird toxisch oder unbrauchbar.
    \item \textbf{3-4, Nebenwirkung}: Unsaubere Arbeit. Ein einsetzendes Unwohlsein verursacht für eine Stunde einen Malus von \textbf{-1 Würfel} auf alle Proben des Konsumenten.
    \item \textbf{5, Invertierung}: Die Substanz wirkt exakt gegenteilig. Ein Heilserum verursacht Wunden, ein Sprengstoff friert die Umgebung ein.
    \item \textbf{6, Katastrophe}: Explosionsartige Ausdehnung. Der Behälter explodiert sofort beim Abfüllen oder Bewegen und verursacht Schaden im direkten Umkreis, zum Beispiel 1W6 Schaden.
\end{itemize}

\begin{pexample}[Beispiel]
Der Überlebenskünstler Garik mischt ein einfaches Elixier (MW 5+, 1 Erfolg) und würfelt mit 4 Würfeln. Das Ergebnis zeigt eine 5 und damit einen Erfolg. Die restlichen drei Würfel landen jedoch auf 1, 1 und 1.

Da Garik drei Einsen, aber nur einen Erfolg gewürfelt hat, gilt die Probe als gepatzt. Die Spielleitung würfelt 1W6 auf der Patzer-Tabelle und erzielt eine 3. Der Trank ist zwar fertig und wirkt, verströmt beim Öffnen aber eine ätzende Gaswolke, die Garik eine Stunde lang einen Malus von \textbf{-1 Würfel} auf alle Proben gibt.
\end{pexample}

\section{Die Goldene Regel}
\index{Goldene Regel}

Alle Modifikatoren, Tabellen und Substitutionsregeln dienen der Atmosphäre. Sie sind optional. Wenn eine Improvisation der Spieler brillant ist, kann die Spielleitung den Malus ignorieren. Wenn ein Patzer erzählerisch keinen Sinn ergibt, kann das Experiment einfach harmlos fehlschlagen. Die gemeinsame Fantasie steht immer über der Tabelle.

\end{multicols}
